\documentclass{situationist-preprint}

\title{The Self Is a Construction, and Real}
\msid{PID-04-a}
\msfamily{PID-04}
\msdiscipline{Personal Identity and the Self}
\mstype{Principal manuscript}
\mskeywords{self; illusion; construction; self-model; no-self; Hume; Metzinger; social ontology}

\begin{document}
\maketitle

\begin{abstract}
Is the self an illusion? A tradition from Hume and the Buddhist doctrine of no-self to self-model theory answers yes: the simple, enduring subject that experience seems to present does not exist. We argue that the self is not an illusion. The self is a construction: it exists in virtue of the activities by which an organism represents, interprets and regulates itself as a single subject and agent. Constructions are real when they figure in explanations that cannot be given without them, and dependence on representational or social activity does not make a thing less real, as MET-20 establishes for social kinds. The self meets this criterion, so it exists. The illusion arguments establish something true but different: the self is misrepresented, in the ordinary self-image, as simple, permanent and independent of the activities that make it. We defend the result against Hume's introspective argument, against Metzinger's claim that self-models represent nothing, and against the objection that the self is a useful fiction.
\end{abstract}

\section{The problem}

Hume looked into himself and found only particular perceptions, never a self that has them, and concluded that a person is nothing but a bundle of perceptions in perpetual flux \citep[1.4.6]{Hume1739}. The Buddhist analysis of the person into five aggregates reaches a similar verdict by a different route: nothing among the aggregates is a self, and there is nothing over and above them \citep{Siderits2003}. Contemporary versions draw on cognitive science. \citet{Metzinger2003} argues that what we call the self is the content of a self-model, a representation the brain constructs of the organism as a whole, which is transparent in the sense that we do not experience it as a representation; nothing in the world corresponds to the enduring subject the model presents, and so no one has ever been or had a self. \citet{Dennett1992} calls the self a center of narrative gravity, an abstraction like a center of mass, posited by the stories a brain spins and not a thing among things.

Realists reply that something answers to the word, whether a thin subject of experience \citep{GStrawson1997} or a person constituted by a self-told story \citep{Schechtman1996}.

We argue that the illusionists are right about how the self is made and wrong about whether it exists: the self is a construction, and constructions are real.

\section{Framework}\label{s:frame}

The question whether X is an illusion has two readings, which the debate does not always keep apart.

\begin{definition}[Illusion; misrepresentation]\label{d:illusion}
Representations as of X are an \emph{existence illusion} if they are systematic and there is nothing that they represent. They are a \emph{nature illusion} if there is something they represent but they represent it as having features it lacks.
\end{definition}

The stick that looks bent in water is a nature illusion; the desert oasis is an existence illusion. The thesis that the self is an illusion, as Hume, Metzinger and Hood state it, is an existence claim: there is no self.

\begin{definition}[Construction]\label{d:construction}
An entity $x$ is a \emph{construction} if $x$ exists at least partly in virtue of activities of representation, interpretation or regulation, including activities directed at $x$ itself.
\end{definition}

Institutions, borders and money are constructions in this sense. The constructing need not be deliberate, and the activities may be subpersonal.

\begin{definition}[The self]\label{d:self}
The \emph{self} of an organism is the unified and ongoing structure of self-representation, self-interpretation and self-regulation by which the organism represents itself as a single subject of experience and agent, ascribes its states and actions to that subject, and governs its conduct by those ascriptions.
\end{definition}

Definition~\ref{d:self} identifies the self by its role. It does not identify the self with the organism, which PID-05 argues is what you are, nor with the person whose stages psychological continuity unifies (PID-01); it identifies the self with a structure the organism sustains, which in ordinary cases is the organism's way of being a subject to itself.

\begin{principle}[Explanatory criterion]\label{p:crit}
A kind exists if it has instances and figures in true explanations that cannot be reformulated, without loss of explanatory power, in terms of other kinds.
\end{principle}

\begin{principle}[Dependence neutrality]\label{p:neutral}
That an entity exists in virtue of representational, interpretive or social activity does not, by itself, make it less real than an entity that does not.
\end{principle}

Principle~\ref{p:crit} is the sufficient condition for existence that MET-20-a uses for social kinds and that underlies realism about natural kinds (MET-27). Principle~\ref{p:neutral} is the general form of the conclusion of MET-20, that the dependence of social kinds on collective attitudes makes them no less real than natural kinds. The argument of MET-20-a does not appeal to anything specific to collective attitudes, and so supports the general principle.

\section{The argument}\label{s:arg}

\begin{lemma}\label{l:constr}
The self is a construction.
\end{lemma}

\begin{proof}
By Definition~\ref{d:self}, the self is a structure of self-representation, self-interpretation and self-regulation. Such a structure exists only while the activities that compose it occur: an organism that ceased to represent itself as a subject, to ascribe its states to that subject and to govern its conduct by the ascriptions would have no self in the defined sense, whatever else it retained. The activities are directed at the structure they compose, since what the organism represents and regulates is itself as a subject. By Definition~\ref{d:construction}, the self is a construction.
\end{proof}

Lemma~\ref{l:constr} concedes to the illusionist what is true in self-model theory: the self is made by the representing. It also agrees with PID-13 that there is no permanent self, since the self persists only as long as its constituent activities. The analogy with institutions is imperfect: an institution, once accepted, persists until it is collectively rejected (MET-20-a), while the self requires the continued occurrence of its activities. The difference concerns persistence, not reality.

\begin{lemma}\label{l:expl}
The self figures in true explanations that cannot be reformulated without it.
\end{lemma}

\begin{proof}
Two classes of explanation illustrate the claim. First, information encoded by relating it to oneself is remembered better than information encoded in other ways, including by attention to its meaning alone \citep{RogersKuiperKirker1977}; the explanation cites the organization of memory around a self-representation, and a description that mentions only the particular items encoded misses the generalization. Second, self-control is explained by an agent's representing a present choice as one of a series of choices of a single agent, so that yielding now is seen as a policy rather than an episode; a description of the episodes, without the representation that unifies them, does not predict which temptations are resisted. In both cases, replacing the structure by a list of the states it organizes loses the generalization, and replacing it by the organism's physiology loses the level at which the regularities hold.
\end{proof}

\begin{theorem}\label{t:main}
The self is not an illusion. It is a construction, and it is real.
\end{theorem}

\begin{proof}
Every normal human organism sustains a structure of the kind described in Definition~\ref{d:self}, so the kind has instances. By Lemma~\ref{l:expl}, it figures in explanations that cannot be reformulated without it. By Principle~\ref{p:crit}, the self exists. By Lemma~\ref{l:constr}, it is a construction, and by Principle~\ref{p:neutral} its being one does not make it less real. Since there is something that self-representations represent, namely the structure they help to compose, they are not an existence illusion in the sense of Definition~\ref{d:illusion}.
\end{proof}

\begin{corollary}\label{c:nature}
The arguments that the self is an illusion establish, at most, a nature illusion: the ordinary self-image represents the self as simple, permanent and independent of the activities that compose it, and it is none of these.
\end{corollary}

\begin{proof}
The self of Theorem~\ref{t:main} is a structure of many activities, so it is not simple; it persists only as long as they do (Lemma~\ref{l:constr}), so it is not permanent; and it exists in virtue of them, so it is not independent of them. To the extent that the ordinary self-image presents it otherwise, it misrepresents the self's nature. By Theorem~\ref{t:main}, there is a self, so the misrepresentation is not an existence illusion.
\end{proof}

Self-representation is partly self-fulfilling: in representing itself as a single agent and regulating its conduct accordingly, the organism makes itself into one. The representation does not represent its own role in building the self, which therefore appears found rather than made. This is the sense in which Metzinger's transparency is correct.

\section{Objections and replies}\label{s:obj}

\begin{objection}[Hume]
Introspection never finds a self, only particular perceptions. If there were a self, attention to one's mind would reveal it.
\end{objection}

\begin{reply}
The objection assumes that the self, if it existed, would be an item among the perceptions, available to inner observation as a sensation is. On Definition~\ref{d:self}, the self is the organization of perceptions and actions under self-ascription, not a further perception. Looking among the perceptions for the self is like looking among the players for the team. Hume's own comparison of the mind to a republic \citep[1.4.6]{Hume1739} is closer to the view defended here than to the denial of a self. The companion manuscript PID-04-b develops the point.
\end{reply}

\begin{objection}[Self-models represent nothing]
The self-model represents a single, enduring subject. Nothing in the brain or the organism is such a subject; there are only processes, among them the self-model. So the self-model has no referent, and the self is an existence illusion \citep{Metzinger2003}.
\end{objection}

\begin{reply}
The objection moves from the claim that nothing has the features the model attributes to the claim that the model represents nothing. That inference is invalid for representations in general: a map that shows a coastline as smoother than it is still represents the coastline. The self-model represents the organism as a unified subject and agent. The organism, through the very activities the model participates in, is unified as a subject and agent to the degree the explanations of Lemma~\ref{l:expl} require. What it is not is unified as a simple substance, which is the content Corollary~\ref{c:nature} classes as misrepresentation.
\end{reply}

\begin{objection}[Useful fiction]
The self is like a center of gravity: an abstraction useful for prediction that is not one of the things in the world \citep{Dennett1992}. Usefulness is not existence.
\end{objection}

\begin{reply}
A center of gravity is an abstractum whose explanatory work is wholly derivative from the distribution of mass, and which does nothing that the distribution does not do. The self is not like that: the self-representation is itself a cause of what it represents, since the organism's regulating its conduct by its self-ascriptions is part of what unifies its conduct. Moreover, on Principle~\ref{p:crit}, a kind that figures ineliminably in true explanations exists; the label \enquote{fiction} adds nothing unless it marks a failure of that criterion. PID-09 reaches the corresponding conclusion for the narrative self: narrative identity is real, though co-authored by others and distinct from numerical identity.
\end{reply}

\section{What the result does not establish}\label{s:not}

The argument does not identify the self with the thing you are, which PID-05 argues is a thinking organism; the self is a structure that organism sustains, and its persistence can diverge from that of the person (PID-01, PID-08). It supports PID-13's conclusion that the self is not permanent. It leaves to PID-09 how much of the self is narrative, to PID-12 how much is constituted by recognition, and to PID-15 how unified it must be.

\section{Conclusion}

The self is made, not found: it is a structure of self-representation, self-interpretation and self-regulation that an organism sustains. Being made does not make it unreal; it figures ineliminably in explanation, and by the criterion that establishes the reality of social kinds, it exists. What the illusionist tradition has shown is that the self is misrepresented as simple, permanent and independent; it has not shown that there is nothing there.

\begin{thebibliography}{99}\small
\bibitem[Dennett(1992)]{Dennett1992} Dennett, D.~C. (1992). The self as a center of narrative gravity. In F.~S. Kessel, P.~M. Cole and D.~L. Johnson (eds.), \emph{Self and Consciousness: Multiple Perspectives} (pp.~103--115). Lawrence Erlbaum.
\bibitem[Hume(1739)]{Hume1739} Hume, D. (1739). \emph{A Treatise of Human Nature}. Cited by book, part and section.
\bibitem[Metzinger(2003)]{Metzinger2003} Metzinger, T. (2003). \emph{Being No One: The Self-Model Theory of Subjectivity}. MIT Press.
\bibitem[Rogers et~al.(1977)]{RogersKuiperKirker1977} Rogers, T.~B., Kuiper, N.~A., and Kirker, W.~S. (1977). Self-reference and the encoding of personal information. \emph{Journal of Personality and Social Psychology}, 35(9), 677--688.
\bibitem[Schechtman(1996)]{Schechtman1996} Schechtman, M. (1996). \emph{The Constitution of Selves}. Cornell University Press.
\bibitem[Siderits(2003)]{Siderits2003} Siderits, M. (2003). \emph{Personal Identity and Buddhist Philosophy: Empty Persons}. Ashgate.
\bibitem[Strawson(1997)]{GStrawson1997} Strawson, G. (1997). The self. \emph{Journal of Consciousness Studies}, 4(5--6), 405--428.
\end{thebibliography}

\end{document}
